\documentclass[11pt]{article}
\usepackage[margin=1in]{geometry}
\usepackage{amsmath,amssymb,amsthm}
\usepackage{hyperref}
\usepackage{enumitem}
\usepackage{xurl}
\let\oldus\_
\renewcommand{\_}{\oldus\allowbreak}

\newtheorem{theorem}{Theorem}
\newtheorem{lemma}[theorem]{Lemma}
\theoremstyle{definition}
\newtheorem{definition}[theorem]{Definition}

\newcommand{\Z}{\mathbb{Z}}
\newcommand{\R}{\mathbb{R}}
\newcommand{\N}{\mathbb{N}}

\title{A disproof of Conjecture 00000007794\\[2pt]
\large Reeve tetrahedra break log-concavity of Ehrhart coefficients}
\date{}

\begin{document}
\sloppy
\maketitle

\section{The conjecture}

Conjecture 00000007794 reads (English version):

\begin{quote}
Definition: Log-concavity of Ehrhart coefficients: $L(P,k)$ are the coefficients of the Ehrhart series of
a lattice polytope $P$. Conjecture: The log-concavification satisfies $L(P,k)^2 \ge L(P,k-1)
L(P,k+1)(1 + c/k^2)$ for universal $c$, i.e. log-concavity with explicit stability margin; equality
holds at the simplex; and the real-rootedness of the $h^*$-polynomial has an independent
characterization with all roots of $h^*(P)$ in the annulus $\{z: r \le |z| \le 1\}$ for explicit lower
bound $r$.
\end{quote}

The Chinese version writes the Ehrhart series explicitly as $\sum L(P,k)t^k$ (``ge duobaoxing P de
Ehrhart jishu $\Sigma$ L(P,k)t\^{}k de xishu L(P,k)'': ``the coefficients $L(P,k)$ of the Ehrhart series
$\sum L(P,k)t^k$ of the lattice polytope $P$''), and the margin as ``c pushi'' (``$c$ universal'').

\textbf{Result.} The first clause is false for every constant $c > -1$, which covers every positive
``stability margin'' and plain log-concavity ($c = 0$). The Reeve tetrahedra violate it at $k = 1$.
Since the conjecture is a conjunction, it is false. Everything is formalized in Lean~4 with Mathlib
(Section~\ref{sec:lean}).

\section{Definitions, reading and conventions}

\begin{definition}
A \emph{lattice polytope} is the convex hull of a nonempty finite set of lattice points. For $k \in \N =
\{0,1,2,\dots\}$, $L(P,k) = \#(kP \cap \Z^d)$; these are the coefficients of the Ehrhart series $\sum_{k
\ge 0} L(P,k) t^k$. Source: Wikipedia, \emph{Ehrhart polynomial}
(\url{https://en.wikipedia.org/wiki/Ehrhart_polynomial}): $L(P,t)$ is ``the number of lattice points
contained in the polytope tP'', and the Ehrhart series is $\mathrm{Ehr}_P(z) = \sum_{t \ge 0} L(P,t)
z^t$.
\end{definition}

\begin{definition}[Reeve tetrahedron]
For an integer $r \ge 1$, $T_r = \mathrm{conv}\{(0,0,0), (1,0,0), (0,1,0), (1,1,r)\} \subset \R^3$.
Source: Wikipedia, \emph{Reeve tetrahedra} (\url{https://en.wikipedia.org/wiki/Reeve_tetrahedra}):
the vertices are as above, ``where r is a positive integer''.
\end{definition}

\textbf{Reading.} The first clause says: there is a constant $c$, independent of $P$ and $k$, such that
for every lattice polytope $P$ and every $k \ge 1$,
\[ L(P,k)^2 \ \ge\ L(P,k-1)\, L(P,k+1)\, (1 + c/k^2). \]
The range $k \ge 1$ is the one where $c/k^2$ and $L(P,k-1)$ are both defined. We refute this for every
real $c > -1$, in dimension $3$. A statement for all dimensions in particular covers $d = 3$.

\textbf{Conventions.} (i) $L(P,0) = \#(0 \cdot P \cap \Z^3) = \#\{0\} = 1$; this is computed from the
definition (Lean: \texttt{ehrhart\_zero}), and it agrees with the constant term $1$ of the Ehrhart
series. (ii) $kP$ is the pointwise dilation $\{k x : x \in P\}$. (iii) Counting uses Mathlib's
\texttt{Set.ncard}, which is $0$ on infinite sets. Our lower bound for $L(T_r,2)$ therefore includes a
proof that the lattice points of $2T_r$ form a finite set. (iv) $k - 1$ is natural subtraction, which is
exact for $k \ge 1$. (v) Full-dimensionality is not part of the definition of a lattice polytope;
anyway $T_r$ is a genuine tetrahedron of volume $r/6$.

\section{Proof}

\begin{lemma}[facet inequalities]\label{lem:facets}
For $s \ge 0$, every $x \in sT_r$ satisfies $0 \le x_3$, $x_3 \le r x_1$, $x_3 \le r x_2$ and $r(x_1 +
x_2) \le s r + x_3$.
\end{lemma}
\begin{proof}
For $s = 1$, the four vertices satisfy these inequalities. The set they define is an intersection of
half-spaces, hence convex, so it contains the convex hull $T_r$. Multiply by $s \ge 0$.
\end{proof}

\begin{lemma}\label{lem:L1}
$L(T_r, 1) \le 4$ for $r \ge 1$: every lattice point of $T_r$ is a vertex.
\end{lemma}
\begin{proof}
Let $(x,y,z) \in \Z^3 \cap T_r$. If $z = 0$: $r x \ge 0$ and $r y \ge 0$ give $x, y \ge 0$, and $r(x+y)
\le r$ gives $x + y \le 1$; so $(x,y) \in \{(0,0),(1,0),(0,1)\}$. If $z > 0$: $rx \ge z > 0$ gives $x \ge
1$, and likewise $y \ge 1$. Then $2r \le r(x+y) \le r + z$ gives $z \ge r$, and $2z \le r(x+y) \le r + z$
gives $z \le r$. So $z = r$, and $r(x+y) \le 2r$ forces $x = y = 1$.
\end{proof}

\begin{lemma}\label{lem:L2}
$L(T_r, 2) \ge r + 1$ for $r \ge 1$.
\end{lemma}
\begin{proof}
\emph{Finiteness.} By Lemma~\ref{lem:facets} with $s = 2$, a lattice point $(x,y,z)$ of $2T_r$ has $x, y
\ge 0$, $2z \le r(x+y) \le 2r + z$ (so $z \le 2r$), and $r(x+y) \le 4r$ (so $x + y \le 4$). It therefore
lies in the box $[0,4] \times [0,4] \times [0,2r]$.
\emph{Points.} For $0 \le z \le r$, put $w = z/(2r) \in [0, \tfrac12]$. The convex combination with
weights $w, \tfrac12 - w, \tfrac12 - w, w$ of the vertices $(0,0,0), (1,0,0), (0,1,0), (1,1,r)$ is
$(\tfrac12, \tfrac12, \tfrac z2) \in T_r$. So $(1,1,z) \in 2T_r$. These are $r + 1$ distinct lattice points.
\end{proof}

\begin{theorem}\label{thm:main}
Let $c > -1$, and let $r \ge 1$ satisfy $(r+1)(1+c) > 16$, for example $r = \lceil 16/(1+c) \rceil +
1$. Then
\[ L(T_r,1)^2 \ <\ L(T_r,0)\, L(T_r,2)\, (1 + c/1^2). \]
Hence no $c > -1$ satisfies the first clause. In particular $T_{16}$ violates plain log-concavity:
$L(T_{16},1)^2 \le 16 < 17 \le L(T_{16},0) L(T_{16},2)$.
\end{theorem}
\begin{proof}
$L(T_r,1)^2 \le 16 < (r+1)(1+c) \le 1 \cdot L(T_r,2) \cdot (1+c)$, by Lemmas~\ref{lem:L1} and~\ref{lem:L2},
$L(T_r,0) = 1$ and $1 + c > 0$.
\end{proof}

\noindent\textbf{Cross-check, not used in the proof.} Wikipedia (\emph{Reeve tetrahedra}) states:
``The Ehrhart polynomial of the Reeve tetrahedron Tr of height r is'' $\frac r6 t^3 + t^2 + (2 - \frac r6)t
+ 1$. This gives $L = 1, 4, r+9, 4r+16, \dots$. A brute-force count confirmed these values for $r \in
\{1,2,8,9,13\}$ and $k \le 3$. With the exact value $r + 9$, plain log-concavity already fails at $r =
8$ ($16 < 17$).

\section{Lean formalization}\label{sec:lean}

File \texttt{lean/Conjecture7794/Basic.lean} (251 lines, namespace \texttt{C7794}; only import
\texttt{Mathlib}). Definitions: \texttt{E3 = Fin 3 -> R}, \texttt{toR}, \texttt{IsLatticePolytope}
(\texttt{convexHull} of the image of a nonempty \texttt{Finset} of integer points),
\texttt{latticePts}, \texttt{ehrhart P k = (latticePts ((k:R) . P)).ncard}, \texttt{vZ}, \texttt{V},
\texttt{reeve}, \texttt{H}. Lemmas: \texttt{reeve\_isLatticePolytope}, \texttt{convex\_H},
\texttt{reeve\_subset\_H}, \texttt{smul\_reeve\_subset} (Lemma~\ref{lem:facets}); \texttt{ehrhart\_zero};
\texttt{latticePts\_one}, \texttt{ehrhart\_one\_le} (Lemma~\ref{lem:L1}); \texttt{latticePts\_two\_finite},
\texttt{mem\_two\_reeve}, \texttt{ehrhart\_two\_ge} (Lemma~\ref{lem:L2}); \texttt{reeve\_violates}
(Theorem~\ref{thm:main}). Main theorem (ASCII rendering):

\begin{verbatim}
theorem conjecture_7794_false :
    not (exists c : R, -1 < c and forall P : Set E3, IsLatticePolytope P ->
        forall k : N, 1 <= k ->
          (ehrhart P (k - 1) : R) * ehrhart P (k + 1) * (1 + c / (k : R) ^ 2)
            <= (ehrhart P k : R) ^ 2) and
    (forall c : R, -1 < c -> exists r : N, 1 <= r and
        IsLatticePolytope (reeve r) and
        (ehrhart (reeve r) 1 : R) ^ 2 <
          (ehrhart (reeve r) 0 : R) * ehrhart (reeve r) 2
            * (1 + c / ((1 : N) : R) ^ 2)) and
    (ehrhart (reeve 16) 1) ^ 2
      < ehrhart (reeve 16) 0 * ehrhart (reeve 16) 2
\end{verbatim}

\textbf{Axiom audit.} \texttt{\#print axioms C7794.conjecture\_7794\_false} reports
\texttt{[propext, Classical.choice, Quot.sound]}. There is no \texttt{sorry} and no
\texttt{native\_decide}. The file compiles under \texttt{lake build} and with
\texttt{-DwarningAsError=true}.

\section{Scope}

\begin{itemize}
\item Refuted: the first clause for every universal $c > -1$ (all positive margins and $c = 0$), with $k$
ranging over $k \ge 1$, in dimension 3. The planar quadrilaterals suggested by an automated filter are
\emph{not} used: by Pick's theorem, lattice polygons satisfy plain log-concavity at $k = 1$.
\item Not refuted: $c \le -1$. Then $1 + c/k^2 \le 0$ at $k = 1$, so the inequality is not a stability
margin at all. Also not refuted: a reading restricted to $k \ge 2$. Reeve tetrahedra satisfy plain
log-concavity at $k = 2$ ($(r+9)^2 > 4(4r+16)$), so they say nothing about small $c > 0$ there.
\item The clauses ``equality holds at the simplex'' and the $h^*$ annulus are not addressed. Note that
$T_r$ is itself a lattice simplex.
\end{itemize}

\end{document}
